\section*{अध्याय \ref{ch:g6:producing-our-food} --- अपना भोजन पैदा करना}

\begin{solution}{exo:g6:producing-our-food:1}
\omterm{def:g6:producing-our-food:farming}{फ़सल उगाना} \omterm{def:g5:food-webs:roles}{उत्पादकों} को सँभालता है --- यानी वे \omterm{def:g1:plants-around-us:plant}{पौधे} जो अपने \omterm{def:g6:plants-are-producers:matter}{कार्बनिक
पदार्थ} के लिए उगाए जाते हैं। \omterm{def:g6:producing-our-food:farming}{पशुपालन} \omterm{def:g5:food-webs:roles}{उपभोक्ताओं} को सँभालता है --- यानी
वे \omterm{def:g1:animals-around-us:animal}{जंतु} जो अपने उत्पादन के लिए पाले जाते हैं: \omterm{def:g2:animals-grow-up:born}{दूध}, माँस, \omterm{def:g2:animals-grow-up:born}{अंडे}, ऊन।
\end{solution}

\begin{solution}{exo:g6:producing-our-food:2}
अनाज: गेहूँ के \omterm{def:g1:plants-around-us:plant}{पौधे} का \omterm{def:g2:from-seed-to-plant:seed}{बीज} (उसका बँधा हुआ भंडार)। आलू: आलू के \omterm{def:g1:plants-around-us:plant}{पौधे} का
\omterm{ex:g4:plants-without-seeds:bulbtuber}{कंद}। सेब: सेब के \omterm{def:g1:plants-around-us:tree}{पेड़} का \omterm{prop:g3:flowers-fruits-seeds:transform}{फल}। \omterm{def:g2:animals-grow-up:born}{दूध}: गाय की माँ-देखभाल का उत्पाद। \omterm{def:g2:animals-grow-up:born}{अंडा}:
मुर्ग़ी का बँधा हुआ पालनाघर।
\end{solution}

\begin{solution}{exo:g6:producing-our-food:3}
\omterm{def:g6:producing-our-food:fermentation}{सूक्ष्मजीवों} के पोषण लेने से किसी कच्चे माल का बदल जाना। रोटी (कच्चा
माल: गूँधा हुआ आटा, मज़दूर: \omterm{ex:g6:producing-our-food:bread}{खमीर}) और दही (कच्चा माल: \omterm{def:g2:animals-grow-up:born}{दूध}, मज़दूर: खट्टा
करने वाले सूक्ष्मजीव); पनीर और सिरका भी चलेंगे।
\end{solution}

\begin{solution}{exo:g6:producing-our-food:4}
\omterm{ex:g6:producing-our-food:bread}{खमीर} --- यानी एककोशिकीय \omterm{def:g6:decomposers-and-soil:decomposers}{कवक} --- आटे की शक्कर पर पलता है और एक गैस
छोड़ता है; लचीला आटा उस गैस को बुलबुलों के रूप में फँसाकर \omterm{def:g3:flowers-fruits-seeds:parts}{फूल} जाता है।
ओवन की गरमी \omterm{ex:g6:producing-our-food:bread}{खमीर} का काम ख़त्म कर देती है और फूला हुआ आटा छेदों समेत जमा
देती है।
\end{solution}

\begin{solution}{exo:g6:producing-our-food:5}
वह अम्ल, जिसे \omterm{def:g2:animals-grow-up:born}{दूध} की शक्कर पर पलते सूक्ष्मजीव छोड़ते हैं: वही अम्ल \omterm{def:g2:animals-grow-up:born}{दूध}
को गाढ़ा करके जमा देता है, और वह ताज़ा खटास ख़ुद वही अम्ल है।
\end{solution}

\begin{solution}{exo:g6:producing-our-food:6}
नमक लगाना, ठंडा रखना (फ़्रिज), सुखाना --- और पकाना भी: हर एक अनचाहे
\omterm{def:g6:producing-our-food:fermentation}{सूक्ष्मजीवों} से काम की कोई परिस्थिति (नमी, गरमाहट) छीन लेता है, ठीक
वैसे ही जैसे ठंड खाद के ढेर को धीमा कर देती है।
\end{solution}

\begin{solution}{exo:g6:producing-our-food:7}
पानी: किसान बारिश पर निर्भर रहता है (और कभी-कभी ख़ुद पानी लाता है)।
\omterm{def:g6:exploring-our-environment:conditions}{प्रकाश}: खुले खेतों में बुआई, और इतनी दूरी पर कि हर \omterm{def:g1:plants-around-us:plant}{पौधे} को अपना हिस्सा
मिले। गरमाहट: गरम मौसम के हिसाब से बुआई का समय। हवा: खुली ज़मीन, ढीली
की हुई \omterm{def:g6:decomposers-and-soil:soil}{मिट्टी}। \omterm{prop:g4:what-plants-need:needs}{खनिज लवण}: खाद और \omterm{def:g6:producing-our-food:farming}{फ़सल} के बाद भरपाई।
\end{solution}

\begin{solution}{exo:g6:producing-our-food:8}
(1) जीव और अवस्था: गाय का \omterm{def:g2:animals-grow-up:born}{दूध} --- यानी \omterm{def:g5:food-webs:roles}{उपभोक्ता} का उत्पाद। (2) चक्र
पर जगह: \omterm{def:g5:food-webs:roles}{उपभोक्ता} का उत्पाद, \omterm{def:g6:producing-our-food:fermentation}{सूक्ष्मजीवों} से पूरा किया हुआ। (3) कच्चा
माल \omterm{def:g2:animals-grow-up:born}{दूध}; मज़दूर उसे खट्टा करके दही में जमा देते हैं, और फिर दबाए हुए
पनीर को महीनों पकाते हैं --- छेद, ऊपरी परत और महक उनका रोज़नामचा हैं।
(4) रास्ता: \omterm{def:g2:animals-grow-up:born}{दूध} गाय से, गाय घास से, घास \omterm{def:g6:exploring-our-environment:conditions}{प्रकाश} से।
\end{solution}

\begin{solution}{exo:g6:producing-our-food:9}
सूक्ष्मजीव अपनी \omterm{def:g6:exploring-our-environment:conditions}{परिस्थितियों} वाले जीवित मज़दूर हैं: ठंड उनका खाना लगभग
रोक देती है। फ़्रिज में यह भोजन की रक्षा करता है (बिगाड़ने वाले बेकार
बैठ जाते हैं); और सर्दी के ढेर में यह खाद बनने में देर कर देता है
(\omterm{def:g6:decomposers-and-soil:decomposers}{अपघटक} बेकार बैठ जाते हैं)। सिद्धांत यह: \omterm{prop:g6:decomposers-and-soil:decomposition}{अपघटन} और बिगड़ना उसी रफ़्तार
से चलते हैं जिसकी \omterm{def:g6:exploring-our-environment:conditions}{परिस्थितियाँ} इजाज़त दें।
\end{solution}

\begin{solution}{exo:g6:producing-our-food:10}
``सूक्ष्मजीव'' नाम \omterm{def:g1:the-five-senses:sense}{आँखों} से न दिखने वाले अलग-अलग सजीवों की एक विशाल
भीड़ का है। कुछ, ग़लत जगह पहुँचकर, बीमारी करते हैं --- हाथ धोना उन्हीं
को निशाना बनाता है; और दूसरे, चुने और क़ाबू में रखी \omterm{def:g6:exploring-our-environment:conditions}{परिस्थितियों} में
सँभाले जाकर, रोटी फुलाते और \omterm{def:g2:animals-grow-up:born}{दूध} जमाते हैं। दुश्मन या कर्मचारी होना
सूक्ष्मजीव की प्रकृति नहीं, बल्कि \omterm{def:g5:biodiversity:species}{प्रजाति}, जगह और प्रबंध के मेल की बात
है।
\end{solution}

\begin{solution}{exo:g6:producing-our-food:11}
हैम एक सूअर से होकर आया: उसके वज़न ने अपने वज़न से कई गुना चारा माँगा,
और उसके ऊपर सूअर का जला हुआ हिस्सा, ज़मीन तथा पानी भी --- यानी
\omterm{def:g5:food-webs:roles}{उपभोक्ता} मंज़िल का चुंगी-कर। रोटी का गेहूँ \omterm{def:g6:exploring-our-environment:conditions}{प्रकाश} से एक ही मंज़िल दूर
था। सैंडविच बरबाद करना दोनों बहीखाते बरबाद करता है; और हैम का बहीखाता
ज़्यादा लंबा है।
\end{solution}

\begin{solution}{exo:g6:producing-our-food:12}
वे मज़दूरों की \emph{\omterm{def:g6:exploring-our-environment:conditions}{परिस्थितियाँ}} सँभालते थे: आटे के कोने की गरमाहट,
जामन को सँभालकर खिलाते रहना, और पीढ़ियों की आज़माइश से सधे नमक तथा समय।
उनकी तरकीबें इसलिए काम करती थीं कि \omterm{def:g6:exploring-our-environment:conditions}{परिस्थितियाँ} ठीक रखने से सही \omterm{def:g6:exploring-our-environment:components}{सजीव}
फलते-फूलते रहते हैं --- चाहे किसी को पता हो या न हो कि वे \omterm{def:g6:exploring-our-environment:components}{सजीव} वहाँ
हैं। उन्होंने \omterm{def:g6:producing-our-food:fermentation}{सूक्ष्मजीवों} की खेती \omterm{def:g1:the-five-senses:sense}{आँखें} मूँदकर की, और अच्छी की।
\end{solution}

\begin{solution}{pb:g6:producing-our-food:1}
\textbf{1.} हर दाना एक \omterm{def:g2:from-seed-to-plant:seed}{बीज} है: \omterm{def:g3:animals-through-seasons:active}{भोजन-भंडार} से भरा गेहूँ का एक नन्हा
\omterm{def:g1:plants-around-us:plant}{पौधा}, जिसे जनक ने उसी पौधशिशु के पहले दिनों के लिए बाँधा था --- हमारे
लिए नहीं, चाहे हम उसे मोड़ लेते हों।

\textbf{2.} ख़ुद गेहूँ के \omterm{def:g1:plants-around-us:plant}{पौधों} की हरी \omterm{def:g1:plants-around-us:parts}{पत्तियों} में, \omterm{def:g6:decomposers-and-soil:soil}{मिट्टी} के पानी तथा
\omterm{prop:g4:what-plants-need:needs}{खनिज लवणों} और हवा की कार्बन डाइऑक्साइड से, और \omterm{prop:g3:balanced-diet:jobs}{ऊर्जा} के रूप में \omterm{def:g6:exploring-our-environment:conditions}{प्रकाश}
से।

\textbf{3.} बोया गया दाना अपना \omterm{def:g2:born-grow-die:cycle}{जीवन चक्र} चलाता है: \omterm{def:g2:from-seed-to-plant:germination}{अंकुरण} से \omterm{def:g1:plants-around-us:plant}{पौधा} और
फिर नया अनाज। बेचा और पीसा गया दाना अपना \omterm{def:g2:born-grow-die:cycle}{जीवन चक्र} वहीं ख़त्म कर देता
है और उसका \omterm{def:g3:animals-through-seasons:active}{भोजन-भंडार} हमारा हो जाता है --- यानी पौधशिशु के लिए बँधा
भोजन, किसी व्यक्ति का खाया हुआ।

\textbf{4.} जो उधार लो वह लौटाओ: खाद गलकर वे \omterm{prop:g4:what-plants-need:needs}{खनिज लवण} फिर से भर देती है
जिन्हें \omterm{def:g6:producing-our-food:farming}{फ़सल} उठा ले गई थी --- यानी वह \omterm{prop:g4:what-plants-need:needs}{खनिज लवणों} की ज़रूरत की सेवा करती
है।

\textbf{5.} क्योंकि \omterm{ex:g6:producing-our-food:bread}{खमीर} एक जीवित मज़दूर है: गरम में वह पोषण लेता और
तेज़ काम करता है; ठंड में बेकार बैठ जाता है --- यानी खाद के दल का वही
नियम, बेकरी के पैमाने पर।

\textbf{6.} उसने आटे की शक्कर पर पोषण लिया और गैस छोड़ी; और लचीले आटे
ने उस गैस को बुलबुलों के रूप में फँसा लिया, और बुलबुलों ने आटा दुगुना
कर दिया।

\textbf{7.} ठंडी खिड़की ने \omterm{ex:g6:producing-our-food:bread}{खमीर} को बेकार बैठा दिया --- मुश्किल से पोषण,
मुश्किल से गैस। दोनों टोलियों ने गरमाहट की एक अनजाने में हुई न्यायी
जाँच चला दी: वही आटा, वही दिन, एक ही परिस्थिति अलग, और तुलना के लिए दो
फूलन।

\textbf{8.} गरमी ने \omterm{ex:g6:producing-our-food:bread}{खमीर} मार दिया और पाली ख़त्म कर दी; और जमे हुए आटे
में बुलबुलों के छेद बचे रह जाते हैं --- यानी कार्यबल के दस्तख़त, वहीं
पके हुए।

\textbf{9.} टुकड़ा रोटी से, रोटी गूँधे हुए आटे से, आटा पिसे अनाज से,
अनाज गेहूँ के \omterm{def:g1:plants-around-us:plant}{पौधे} की \omterm{def:g1:plants-around-us:parts}{पत्तियों} से --- जहाँ वह खनिज सामग्री से, \omterm{def:g6:exploring-our-environment:conditions}{प्रकाश}
के साथ बनाया गया था: रास्ता जुलाई के खेत की धूप पर ख़त्म होता है।

\textbf{10.} टुकड़ा: एक मंज़िल (\omterm{def:g5:food-webs:roles}{उत्पादक} से मेज़ तक, \omterm{def:g6:producing-our-food:fermentation}{सूक्ष्मजीवों} की
आख़िरी छुअन के साथ)। मक्खन: दो मंज़िलें (घास से गाय, गाय से मलाई)। हैम:
दो मंज़िलें और सबसे बड़ा चुंगी-कर --- सूअर ने अपने चारे का ज़्यादातर
पदार्थ हैम बनने से पहले ही जला और छोड़ दिया था। टुकड़ा पदार्थ में सबसे
सस्ता है; हैम सबसे महँगा।

\textbf{11.} खाद के \omterm{def:g6:decomposers-and-soil:decomposers}{अपघटक} --- केंचुए, कठ-जूएँ, \omterm{def:g6:decomposers-and-soil:decomposers}{कवक}, और न दिखने वाले दल
--- आटे के \omterm{def:g6:plants-are-producers:matter}{कार्बनिक पदार्थ} को खोलकर \omterm{def:g6:decomposers-and-soil:soil}{ह्यूमस} और \omterm{prop:g4:what-plants-need:needs}{खनिज लवण} बना देते हैं:
यानी भोजन से अगले साल की \omterm{def:g6:decomposers-and-soil:soil}{मिट्टी} की उर्वरता तक, चक्र का एक मोड़ और।

\textbf{12.} जैसे: ``हमारी रोटी को एक \omterm{def:g5:food-webs:roles}{उत्पादक} ने \omterm{def:g6:plants-are-producers:matter}{खनिज पदार्थ} और \omterm{def:g6:exploring-our-environment:conditions}{प्रकाश}
से बनाया, \omterm{def:g6:producing-our-food:fermentation}{सूक्ष्मजीवों} के कार्यबल ने फुलाया, और उसमें कोई \omterm{def:g5:food-webs:roles}{उपभोक्ता}
मंज़िल लगी ही नहीं --- हालाँकि बग़ल में बिकते हैम के सैंडविचों में एक
लगी थी।''
\end{solution}
